\documentclass[11pt,a4paper]{article}
\usepackage[T1]{fontenc}
\usepackage{lmodern}
\usepackage[margin=27mm]{geometry}
\usepackage{amsmath,amssymb,amsthm,mathtools}
\usepackage{microtype}
\usepackage{booktabs,array}
\usepackage{enumitem}
\usepackage[hidelinks]{hyperref}
\setlength{\parskip}{5pt plus 1pt}
\setlength{\parindent}{0pt}
\setlength{\emergencystretch}{3em}
\setlist{itemsep=2pt,topsep=3pt}
\raggedbottom
\newtheorem{theorem}{Theorem}[section]
\newtheorem{lemma}[theorem]{Lemma}
\newtheorem{corollary}[theorem]{Corollary}
\newtheorem{proposition}[theorem]{Proposition}
\theoremstyle{definition}
\newtheorem{definition}[theorem]{Definition}
\newtheorem{example}[theorem]{Example}
\theoremstyle{remark}
\newtheorem{remark}[theorem]{Remark}
\newcommand{\Z}{\mathbb{Z}}
\newcommand{\N}{\mathbb{N}}
\newcommand{\Q}{\mathbb{Q}}
\newcommand{\F}{\mathbb{F}}
\newcommand{\ps}[1]{#1[[x]]}
\newcommand{\coeff}[2]{[x^{#1}]\,#2}
\newcommand{\hl}[1]{\textbf{#1}}
\newcommand{\para}[1]{\par\medskip\noindent\textbf{#1}\enspace}
\hypersetup{pdftitle={Twice the squares, not the squares: a disproof of the conjecture in OEIS A380680, and the corrected statement},pdfauthor={Dev D. Goyal},pdfsubject={OEIS A380680}}
\title{Twice the squares, not the squares: remainders modulo 4 in a theta-type series\\[6pt]\large A disproof of the conjecture in OEIS A380680, and the corrected statement}
\author{Dev D. Goyal}
\date{October 2026}
\begin{document}
\maketitle
\begin{abstract}
The OEIS entry A380680, contributed by Paul D.~Hanna, lists the coefficients of the power series
$A(x)=\sum_{n=-\infty}^{\infty}x^{n^2}\bigl((2-x^n)/(1-2x^n)\bigr)^n=1+4x+6x^2+12x^3+32x^4+\cdots$, a variant of the
classical theta series $\theta(x)=\sum_nx^{n^2}$. The entry conjectures that, for $n>0$, the coefficient $a(n)$ leaves
remainder $2$ on division by $4$ exactly when $n$ is a perfect square, and is divisible by $4$ otherwise. The published
terms $a(1)=4$ and $a(2)=6$ already contradict the conjecture. We prove the correct statement: for $n>0$, $a(n)$ leaves remainder $2$ exactly when $n$ is twice a perfect square ($n=2,8,18,32,\dots$),
and $a(n)$ is divisible by $4$ otherwise. Equivalently, $A(x)\equiv\theta(x^2)\pmod 4$. The results have also been checked by computer and proved formally in the Lean proof assistant.
\end{abstract}
\begin{center}\begin{minipage}{0.88\linewidth}\small
\textbf{How this paper was produced.} This paper reports the result of an AI research effort directed by
Dev~D.~Goyal. ChatGPT (OpenAI), using the GPT-6 model with Pro reasoning and working autonomously from prompts by
Dev~D.~Goyal, found the proof, wrote its Lean~4 formalization, and wrote a first draft of the paper.
Dev~D.~Goyal chose the problems and asked for proofs, formal verification and revisions, but did not supply proof
ideas. Claude Opus~5.5, used with Claude Code, then checked the work independently: it rebuilt the formal proofs from
source and confirmed that they rely only on Lean's three standard axioms, recomputed the sequence from its definition
compared it with the published terms and tested every claim numerically, and read every step of the proofs. It also wrote the present
version of this paper.
Dev~D.~Goyal has reviewed the proofs and is responsible for them; this review is not an independent expert validation.
The paper has not been peer reviewed.
The section ``About this paper'' at the end gives the details.
\end{minipage}\end{center}

\section{Introduction}

\subsection{The sequence}\label{sec:sequence}
The \emph{On-Line Encyclopedia of Integer Sequences} (OEIS)~\cite{oeisdb} is a large public database of integer
sequences. Many
entries also record \emph{conjectures}: patterns that the contributor observed in the data but did not prove.

This paper is about entry A380680~\cite{oeis}, contributed by Paul~D.~Hanna in February 2025. It reads:
\begin{quote}\small\ttfamily\raggedright G.f. A(x) = Sum\_\{n=-oo..+oo\} x\^{}(n\^{}2) * (2 - x\^{}n)\^{}n/(1 - 2*x\^{}n)\^{}n.\par\smallskip The term x\^{}(n\^{}2) * (2 - x\^{}n)\^{}n/(1 - 2*x\^{}n)\^{}n is invariant under sign reversal of n.\par\smallskip\hfill\textrm{--- OEIS A380680, Paul D. Hanna, February 2025}\end{quote}

Here ``G.f.''\ stands for \emph{generating function}: the series $A(x)=\sum_{n\ge0}a(n)\,x^n$ whose coefficients are
the terms of the sequence, and \texttt{Sum\_\{n=-oo..+oo\}} is a sum over all integers $n$. This series is given by an explicit formula: the right-hand side does not involve $A$. Using the symmetry
noted in the entry, the sum over all integers can be folded into a sum over the positive integers
(Lemma~\ref{lem:fold} below with $U=2$), and the definition becomes
\begin{equation}\label{eq:def}
A=1+2\,\Sigma(2),\qquad\text{where}\qquad
\Sigma(2)=\sum_{j\ge1}x^{j^2}\Bigl(\frac{2-x^j}{1-2x^j}\Bigr)^{j}.
\end{equation}
Here $2$ stands for the constant series. This one-sided form is formula~(2) of the entry (not to be confused
with equation~\eqref{eq:fold} of this paper), and it is the form we work with.

The first terms, starting at $n=0$, are 1, 4, 6, 12, 32, 48, 120, 192, 450, 784, 1704, 3072, 6624, 12288, \dots{}

The sum in the definition looks like the classical \emph{theta series} $\theta(x)=\sum_nx^{n^2}=1+2x+2x^4+2x^9+\cdots$
(Section~\ref{sec:background}), except that each term $x^{n^2}$ carries an extra factor. We call such sums
\emph{theta-type sums}. A companion entry by the same author, OEIS A380060~\cite{oeisother}, contributed on
the same day, defines a series $A$ by an implicit equation, with $A$ itself in place of the constant $2$ inside the sum,
and states the same conjecture word for word. For that entry the conjecture is true; it is treated in a companion
paper~\cite{repo}.

\subsection{The question}
The entry contains the following conjecture.
\begin{quote}\small\ttfamily\raggedright Conjecture: for n > 0, a(n) == 2 (mod 4) iff n is a square, and a(n) is divisible by 4 when n is nonsquare.\par\smallskip\hfill\textrm{--- OEIS A380680, Paul D. Hanna, February 2025}\end{quote}

(In the OEIS notation, \texttt{==} denotes congruence.) So, for $n>0$, the conjecture says that the terms with remainder $2$ modulo $4$ are exactly those at the
perfect squares $n=1,4,9,16,\dots$. The table shows the remainders of the first $48$ terms, and marks with a dot the
perfect squares and the numbers of the form $2j^2$.
\begin{center}\small\setlength{\tabcolsep}{3.2pt}
\begin{tabular}{>{\raggedright\arraybackslash}p{6.4em}>{\centering\arraybackslash}p{1.15em}>{\centering\arraybackslash}p{1.15em}>{\centering\arraybackslash}p{1.15em}>{\centering\arraybackslash}p{1.15em}>{\centering\arraybackslash}p{1.15em}>{\centering\arraybackslash}p{1.15em}>{\centering\arraybackslash}p{1.15em}>{\centering\arraybackslash}p{1.15em}>{\centering\arraybackslash}p{1.15em}>{\centering\arraybackslash}p{1.15em}>{\centering\arraybackslash}p{1.15em}>{\centering\arraybackslash}p{1.15em}>{\centering\arraybackslash}p{1.15em}>{\centering\arraybackslash}p{1.15em}>{\centering\arraybackslash}p{1.15em}>{\centering\arraybackslash}p{1.15em}}
\toprule
$n$ & 0 & 1 & \hl{2} & 3 & 4 & 5 & 6 & 7 & \hl{8} & 9 & 10 & 11 & 12 & 13 & 14 & 15 \\ \midrule
$a(n)\bmod 4$ & 1 & 0 & \hl{2} & 0 & 0 & 0 & 0 & 0 & \hl{2} & 0 & 0 & 0 & 0 & 0 & 0 & 0 \\
$n=j^2$ &  & $\bullet$ & \hl{} &  & $\bullet$ &  &  &  & \hl{} & $\bullet$ &  &  &  &  &  &  \\
$n=2j^2$ &  &  & \hl{$\bullet$} &  &  &  &  &  & \hl{$\bullet$} &  &  &  &  &  &  &  \\
\bottomrule
\end{tabular}
\\[4pt]
\begin{tabular}{>{\raggedright\arraybackslash}p{6.4em}>{\centering\arraybackslash}p{1.15em}>{\centering\arraybackslash}p{1.15em}>{\centering\arraybackslash}p{1.15em}>{\centering\arraybackslash}p{1.15em}>{\centering\arraybackslash}p{1.15em}>{\centering\arraybackslash}p{1.15em}>{\centering\arraybackslash}p{1.15em}>{\centering\arraybackslash}p{1.15em}>{\centering\arraybackslash}p{1.15em}>{\centering\arraybackslash}p{1.15em}>{\centering\arraybackslash}p{1.15em}>{\centering\arraybackslash}p{1.15em}>{\centering\arraybackslash}p{1.15em}>{\centering\arraybackslash}p{1.15em}>{\centering\arraybackslash}p{1.15em}>{\centering\arraybackslash}p{1.15em}}
\toprule
$n$ & 16 & 17 & \hl{18} & 19 & 20 & 21 & 22 & 23 & 24 & 25 & 26 & 27 & 28 & 29 & 30 & 31 \\ \midrule
$a(n)\bmod 4$ & 0 & 0 & \hl{2} & 0 & 0 & 0 & 0 & 0 & 0 & 0 & 0 & 0 & 0 & 0 & 0 & 0 \\
$n=j^2$ & $\bullet$ &  & \hl{} &  &  &  &  &  &  & $\bullet$ &  &  &  &  &  &  \\
$n=2j^2$ &  &  & \hl{$\bullet$} &  &  &  &  &  &  &  &  &  &  &  &  &  \\
\bottomrule
\end{tabular}
\\[4pt]
\begin{tabular}{>{\raggedright\arraybackslash}p{6.4em}>{\centering\arraybackslash}p{1.15em}>{\centering\arraybackslash}p{1.15em}>{\centering\arraybackslash}p{1.15em}>{\centering\arraybackslash}p{1.15em}>{\centering\arraybackslash}p{1.15em}>{\centering\arraybackslash}p{1.15em}>{\centering\arraybackslash}p{1.15em}>{\centering\arraybackslash}p{1.15em}>{\centering\arraybackslash}p{1.15em}>{\centering\arraybackslash}p{1.15em}>{\centering\arraybackslash}p{1.15em}>{\centering\arraybackslash}p{1.15em}>{\centering\arraybackslash}p{1.15em}>{\centering\arraybackslash}p{1.15em}>{\centering\arraybackslash}p{1.15em}>{\centering\arraybackslash}p{1.15em}}
\toprule
$n$ & \hl{32} & 33 & 34 & 35 & 36 & 37 & 38 & 39 & 40 & 41 & 42 & 43 & 44 & 45 & 46 & 47 \\ \midrule
$a(n)\bmod 4$ & \hl{2} & 0 & 0 & 0 & 0 & 0 & 0 & 0 & 0 & 0 & 0 & 0 & 0 & 0 & 0 & 0 \\
$n=j^2$ & \hl{} &  &  &  & $\bullet$ &  &  &  &  &  &  &  &  &  &  &  \\
$n=2j^2$ & \hl{$\bullet$} &  &  &  &  &  &  &  &  &  &  &  &  &  &  &  \\
\bottomrule
\end{tabular}
\end{center}

The data contradict the conjecture at once: $a(1)=4$ is divisible by $4$ although $1$ is a square, and $a(2)=6$ leaves
remainder $2$ although $2$ is not a square. Instead, the remainder $2$ occurs exactly at $n=2,8,18,32$, which are
twice the squares $1,4,9,16$. These columns are printed in bold.

\subsection{What this paper proves}
\textbf{The conjecture is false} (Corollary~\ref{cor:false}): already $a(1)=4$ and $a(2)=6$ contradict it.

\textbf{Main result} (Theorem~\ref{thm:main}). \emph{For every $n>0$, $a(n)\equiv2\pmod4$ if $n$ is twice a perfect
square, and $a(n)$ is divisible by $4$ otherwise.} Equivalently,
\[
A(x)\equiv\theta(x^2)=1+2\sum_{j\ge1}x^{2j^2}\pmod 4 .
\]
The key observation is that, modulo $2$, the factor $(2-x^j)/(1-2x^j)$ is $x^j$, because $2$ becomes $0$. So the
$j$-th term $x^{j^2}\bigl((2-x^j)/(1-2x^j)\bigr)^j$ becomes $x^{j^2}\cdot x^{j^2}=x^{2j^2}$: the exponents are doubled.

\subsection{What was known}
The entry records the folded form of the sum (formula~(2) in its list of formulas; see Lemma~\ref{lem:fold}) and two further
rewritings of it. The entry does not note that its own data contradict the conjecture. We are not aware of an earlier proof or disproof of the conjecture.

\subsection{Motivation}
The theta series is a classical object of number theory, and Hanna's sums are deformations of it: each term
$x^{n^2}$ is multiplied by a further factor. One may ask how much of the arithmetic of the theta series survives the
deformation. For this entry and its companion A380060, the sums reduce modulo $4$ to theta series
(Section~\ref{sec:discussion}).
The definitions of the two entries differ in only one symbol, their
conjectures are stated in the same words, and yet the conjecture is true for A380060 and false for A380680. The proofs show exactly where the
difference comes from: replacing $A$ by $2$ inside the sum changes each factor $(A-x^j)/(1-x^jA)$, modulo $2$, from $1$
to $x^j$.

\subsection{How to read this paper}
Section~\ref{sec:background} explains the few notions we need: power series, remainders, the theta series, and how to
make sense of a sum over all integers. Section~\ref{sec:results} states the results precisely, and
Section~\ref{sec:proofs} proves them. Section~\ref{sec:checks} describes how the results were
checked, by computer and by a formal proof in Lean. Section~\ref{sec:discussion} discusses the scope of the result
and a generalization.
\section{Background}\label{sec:background}
This section collects the definitions and facts used later. Readers who know formal power series and congruences can skip it and return when needed.
\subsection{Formal power series}
A \emph{formal power series} is an expression
\[
U(x)=u_0+u_1x+u_2x^2+u_3x^3+\cdots ,
\]
determined by its sequence of coefficients $u_0,u_1,u_2,\dots$; see~\cite{niven,stanley}. No number is ever
substituted for $x$, so questions of convergence do not arise. We write $\coeff{n}{U}=u_n$ for the coefficient of $x^n$ in $U$; two series are equal
when all their coefficients are equal.

Series are added coefficient by coefficient and multiplied like polynomials:
\[
\coeff{n}{(UV)}=u_0v_n+u_1v_{n-1}+\cdots+u_nv_0 .
\]
Each coefficient of a product is a finite sum. For example, $(1-x)(1+x+x^2+x^3+\cdots)=1$, so $1+x+x^2+\cdots$ is
the \emph{multiplicative inverse} of $1-x$, written $\frac{1}{1-x}$. More generally, every series $D$ with constant term
$1$ has a unique multiplicative inverse $\frac1D$, whose coefficients are determined one at a time by $D\cdot\frac1D=1$
using only addition, subtraction and multiplication. Whenever a quotient of power series appears in this paper, its
denominator has constant term $1$, so the quotient is again a power series.

A \emph{commutative ring} is a set with an addition, a subtraction and a multiplication that satisfy the usual rules
of arithmetic; examples are $\Z$, $\Q$ and the integers modulo $m$. For a commutative ring $R$ we write $\ps{R}$ for
the set of power series with coefficients in $R$; for instance, writing $\F_2$ for the integers modulo $2$, $\ps{\F_2}$
consists of the series with coefficients modulo $2$. We write
$U(0)$ for the constant term $u_0$ of $U$, and $U=O(x^d)$ when the coefficients of $1,x,\dots,x^{d-1}$ in $U$ are all
zero. Thus $U-V=O(x^d)$ means that $U$ and $V$ have the same coefficient of $x^n$ for every $n<d$. A series with
constant term $0$ and coefficient of $x$ equal to $1$, that is, a series of the form $x+u_2x^2+u_3x^3+\cdots$, is
called \emph{normalized}.

\subsection{Remainders and congruences}
For an integer $m\ge2$ we write $a\equiv b\pmod m$ (``$a$ is congruent to $b$ modulo $m$'') when $a-b$ is divisible by
$m$, that is, when $a$ and $b$ leave the same remainder after division by $m$. The remainder of $a$ on division by
$m$ is called its \emph{residue} modulo $m$; it is one of $0,1,\dots,m-1$.

Addition and multiplication are compatible with congruences, so the residues $0,1,\dots,m-1$, added and multiplied
modulo $m$, form a ring, written $\Z/m\Z$. When $m=p$ is a prime we also write $\F_p$. In $\F_2=\{0,1\}$ we have $1+1=0$,
so addition and subtraction coincide.

For power series with integer coefficients, $U\equiv V\pmod m$ means that \emph{every} coefficient of $U-V$ is
divisible by $m$. Equivalently, the two series become equal when every coefficient is replaced by its residue; we call
this \emph{reducing} the series modulo $m$. The key fact is the following.

\emph{Reduction modulo $m$ is compatible with all the operations above.} To reduce a sum, a product, a substitution, or
an inverse, we may reduce the ingredients first and then compute in $\Z/m\Z$. This holds because every
coefficient of the result is obtained from finitely many coefficients of the ingredients by adding, subtracting and
multiplying, and these operations are compatible with remainders.

\subsection{The theta series}
The \emph{theta series} is the power series
\[
\theta(x)=\sum_{n=-\infty}^{\infty}x^{n^2}=1+2x+2x^4+2x^9+2x^{16}+\cdots .
\]
The sum runs over all integers $n$: positive, negative and zero. The terms for $n$ and $-n$ are equal, and $n=0$
gives $1$, so
\[
\theta(x)=1+2\sum_{j\ge1}x^{j^2}.
\]
So the coefficient of $x^n$ in $\theta$ is the number of integers whose square is $n$: it is $1$ for $n=0$, it is $2$
when $n$ is a positive perfect square (the two square roots $\pm j$), and it is $0$ otherwise. Replacing $x$ by $x^2$
gives $\theta(x^2)=1+2\sum_{j\ge1}x^{2j^2}=1+2x^2+2x^8+2x^{18}+\cdots$, whose nonzero coefficients sit at twice the
squares.

The theta series goes back to Jacobi. For example, the
coefficient of $x^n$ in $\theta(x)^k$ is the number of ways to write $n$ as a sum of $k$ squares of integers,
counting order and signs; see~\cite{grosswald}.

\subsection{The two-sided sum}\label{sec:fold}
Hanna's definitions are sums over all integers $n$, and for negative $n$ the terms contain negative powers of $x$. To
make sense of them we use \emph{formal Laurent series}: series that may also contain finitely many negative powers of
$x$, such as $x^{-2}+3+x+x^2+\cdots$. They are added and multiplied like power series. With rational coefficients,
every nonzero Laurent series has a multiplicative inverse: write it as $c\,x^{-k}(1+\cdots)$ with $c\ne0$ and invert
each factor (the factor $1+\cdots$ has constant term $1$, so it can be inverted as in the previous subsection).

For a power series $U$ and an integer $j\ge1$ we write
\[
Q_j(U)=\frac{U-x^j}{1-x^jU}.
\]
The denominator has constant term $1$, so $Q_j(U)$ is a power series, and it has integer coefficients if $U$
has.

\begin{lemma}[Folding the sum]\label{lem:fold}
Let $U$ be a power series with rational coefficients and nonzero constant term. For every $j\ge1$, the term with index
$n=-j$ of the sum $\sum_{n}x^{n^2}\bigl((U-x^n)/(1-x^nU)\bigr)^n$ is equal to the term with index $n=j$:
\[
x^{j^2}\Bigl(\frac{U-x^{-j}}{1-x^{-j}U}\Bigr)^{-j}=x^{j^2}\,Q_j(U)^j .
\]
Consequently the two-sided sum is a power series, and
\begin{equation}\label{eq:fold}
\sum_{n=-\infty}^{\infty}x^{n^2}\Bigl(\frac{U-x^n}{1-x^nU}\Bigr)^n=1+2\,\Sigma(U),
\qquad\text{where}\qquad
\Sigma(U)=\sum_{j\ge1}x^{j^2}\,Q_j(U)^j .
\end{equation}
\end{lemma}

\begin{proof}
Both $U-x^{-j}$ and $1-x^{-j}U$ are nonzero Laurent series (the first contains the term $-x^{-j}$, and the second is
not zero because $U$ has a nonzero constant term and so is not $x^j$), so the fraction and its inverse make sense. A negative
exponent means: invert, then raise to the power $j$. So the left side is
\[
x^{j^2}\Bigl(\frac{1-x^{-j}U}{U-x^{-j}}\Bigr)^{j}
=x^{j^2}\Bigl(\frac{x^j-U}{x^jU-1}\Bigr)^{j}
=x^{j^2}\Bigl(\frac{U-x^j}{1-x^jU}\Bigr)^{j},
\]
where we multiplied the numerator and the denominator by $x^j$, and then changed the sign of both. The term with
$n=0$ is $1$. Pairing $n=j$ with $n=-j$ gives \eqref{eq:fold}. Finally, the term with index $j$ is $x^{j^2}$ times a
power series, so only the finitely many $j$ with $j^2\le n$ affect the coefficient of $x^n$; this is why the infinite
sum is a well-defined power series.
\end{proof}

The OEIS entry notes the symmetry of Lemma~\ref{lem:fold} (``invariant under sign reversal of $n$'') and records the
folded form as formula~(2) in its list of formulas. In the one-sided form no negative powers occur, and the only division is by a
series with constant term $1$; so it makes sense with coefficients in any commutative ring, including $\Z/4\Z$.

\subsection{Doubling modulo 4}
\begin{lemma}\label{lem:double}
Let $U$ and $V$ be power series with integer coefficients. If $U\equiv V\pmod 2$, then $2U\equiv2V\pmod 4$.
\end{lemma}

\begin{proof}
Every coefficient of $U-V$ is even, so every coefficient of $2U-2V=2(U-V)$ is divisible by $4$.
\end{proof}
In other words, \emph{twice a series is determined modulo $4$ by the series modulo $2$.} The sums in question are
multiplied by $2$, so it suffices to know them modulo $2$.

\section{Results}\label{sec:results}

\begin{theorem}\label{thm:main}
The series $A$ defined by~\eqref{eq:def} satisfies
\[
A(x)\equiv\theta(x^2)=1+2\sum_{j\ge1}x^{2j^2}=1+2x^2+2x^8+2x^{18}+\cdots\pmod4 .
\]
Consequently, for every $n>0$:
\begin{itemize}
\item $a(n)\equiv2\pmod 4$ if $n=2j^2$ for some integer $j\ge1$, and
\item $a(n)$ is divisible by $4$ otherwise.
\end{itemize}
\end{theorem}

\begin{corollary}\label{cor:false}
\begin{enumerate}[label=(\alph*)]
\item $a(1)=4$ and $a(2)=6$. Hence both parts of the conjecture are false: $1$ is a square but $a(1)$ is divisible by
$4$, and $2$ is not a square but $a(2)$ is not divisible by $4$.
\item More precisely, for $n>0$ the conjecture predicts the wrong remainder exactly when $n$ is a perfect square or
twice a perfect square, and the right remainder (namely $0$) for all other $n$.
\end{enumerate}
\end{corollary}

\section{Proofs}\label{sec:proofs}

\para{The idea in brief.} Modulo $2$ the number $2$ is $0$, so the factor $(2-x^j)/(1-2x^j)$ becomes $-x^j/1$, which is $x^j$ because
$-1\equiv1\pmod 2$. Hence the $j$-th term of $\Sigma(2)$ is $x^{j^2}\cdot x^{j^2}=x^{2j^2}$ modulo $2$. Since $\Sigma(2)$
is multiplied by $2$, knowing it modulo $2$ is enough to know $A$ modulo $4$ (Lemma~\ref{lem:double}).

\begin{proof}[Proof of Theorem~\ref{thm:main}]
Reduction modulo $2$ is compatible with sums, products, powers and inverses, so to reduce
$Q_j(2)=(2-x^j)/(1-2x^j)$ we may first reduce the ingredients. The number $2$ becomes $0$, so
\[
Q_j(2)\equiv\frac{0-x^j}{1-0}=-x^j\equiv x^j\pmod 2,
\]
and hence
\[
x^{j^2}Q_j(2)^j\equiv x^{j^2}\,(x^j)^j=x^{2j^2}\pmod 2 .
\]
Summing over $j$ gives $\Sigma(2)\equiv\sum_{j\ge1}x^{2j^2}\pmod2$. By Lemma~\ref{lem:double},
$2\Sigma(2)\equiv2\sum_{j\ge1}x^{2j^2}\pmod4$, and so
\[
A=1+2\Sigma(2)\equiv1+2\sum_{j\ge1}x^{2j^2}=\theta(x^2)\pmod 4 .
\]
Different $j\ge1$ give different exponents $2j^2$, so for $n>0$ the coefficient of $x^n$ in $\theta(x^2)$ is $2$ when
$n=2j^2$ and $0$ otherwise.
\end{proof}

\begin{proof}[Proof of Corollary~\ref{cor:false}]
(a) The terms of $\Sigma(2)$ with $j\ge2$ start at $x^4$, so up to $x^3$ only the term with $j=1$ matters. Expanding,
\[
\frac{2-x}{1-2x}=(2-x)(1+2x+4x^2+\cdots)=2+3x+6x^2+\cdots,
\]
so $A=1+2x\,(2+3x+6x^2+\cdots)+O(x^4)=1+4x+6x^2+12x^3+O(x^4)$. This matches the entry, and it gives $a(1)=4$ and
$a(2)=6$, which contradict the two parts of the conjecture as stated.

(b) For $n>0$ the conjecture predicts remainder $2$ at the squares and $0$ elsewhere; by Theorem~\ref{thm:main} the
truth is remainder $2$ at the numbers $2j^2$ and $0$ elsewhere. No positive integer is both a square and twice a
square: if $j^2=2k^2$ with $j,k\ge1$, then $\sqrt2=j/k$ would be rational, contradicting the irrationality of $\sqrt2$. So the two
descriptions disagree exactly at the squares and at the numbers $2j^2$, and agree everywhere else.
\end{proof}

\section{Checking the results}\label{sec:checks}

A mathematical proof can contain mistakes, so the results were checked in two independent ways.

\para{Numerical check.}
We evaluated the two-sided sum in the OEIS name directly, with the terms for negative $n$ computed as Laurent
series rather than by folding. This reproduces all 36 terms printed in the entry, including $a(1)=4$ and $a(2)=6$. The
pattern of Theorem~\ref{thm:main} holds for every $0<n<220$.
The computation was done with the PARI/GP computer algebra system, using code that was written separately from the
formal proofs (file \texttt{verification/independent\_checks.gp} in the accompanying repository~\cite{repo}).
A computation can only test finitely many cases, so it is a check on the definitions and on the indexing, not a proof.

\para{Formal proof in Lean.}
Lean~\cite{lean} is a \emph{proof assistant}: a program in which mathematical statements and their proofs are written
in a precise formal language, and every step is checked mechanically down to the basic rules of logic. The theorems of Section~\ref{sec:results} are proved in Lean~4 (version 4.35.0) on top of
Mathlib~\cite{mathlib}, Lean's library of formalized mathematics; the steps proved only on paper are identified below. The statements are collected in the file \texttt{lean/Entries/A380680.lean} of the repository,
in the namespace \texttt{OEIS.A380680}. There the one-sided sum $\Sigma(2)$ is written out coefficient by coefficient (only the indices $j\le n$ can contribute to $x^n$), and the multiplicative inverse of a series with constant term $1$ is written \texttt{unitInv}. The folding of the two-sided sum (Lemma~\ref{lem:fold}) is proved in this paper but not in Lean. The proofs are in the file \texttt{lean/Families/ThetaRational.lean}. Part~(b) of Corollary~\ref{cor:false} is a direct consequence of the theorem and is not stated separately in Lean.
The table lists the results covered by the Lean formalization. The Lean formalization does not cover the additional arguments and observations in Section~\ref{sec:discussion}.
\begin{center}\small\renewcommand{\arraystretch}{1.25}
\begin{tabular}{>{\raggedright\arraybackslash}p{0.52\linewidth}>{\raggedright\arraybackslash}p{0.4\linewidth}}
\toprule
Statement in this paper & Lean theorem(s)\\
\midrule
The defining sum~\eqref{eq:def}, coefficient by coefficient & \texttt{explicit\_\allowbreak{}actual\_\allowbreak{}source} \\
Theorem~\ref{thm:main}: $A\equiv\theta(x^2)\pmod4$ & \texttt{complete\_\allowbreak{}series\_\allowbreak{}mod\_\allowbreak{}four} \\
Theorem~\ref{thm:main}: the two statements for each $n>0$ & \texttt{a380680\_\allowbreak{}corrected\_\allowbreak{}statement} \\
Corollary~\ref{cor:false}(a): $a(1)=4$, $a(2)=6$, and both parts of the conjecture are false & \texttt{a380680\_\allowbreak{}counterexamples}, \texttt{a380680\_\allowbreak{}literal\_\allowbreak{}statement\_\allowbreak{}is\_\allowbreak{}false}, \texttt{a380680\_\allowbreak{}complement\_\allowbreak{}is\_\allowbreak{}false} \\
\bottomrule
\end{tabular}
\end{center}
The Lean project was rebuilt from source, and the axioms on which each of these theorems depends were listed with
Lean's \texttt{\#print axioms} command. Every one depends only on the three standard axioms of Lean's logic
(\texttt{propext}, \texttt{Classical.choice}, \texttt{Quot.sound}); no step is left unproved. A formal proof certifies
exactly the statement that Lean checks. Whether that statement says what the OEIS entry means is a separate question,
so the Lean definitions were also compared with the definition in the OEIS entry
(Section~\ref{sec:sequence}). The principal theorems are stated once more, using only Mathlib's definitions, in the
file \texttt{lean/Challenges/Theta.lean}, and Lean's Comparator tool confirms that the
proved theorems have exactly these statements.

\section{Discussion}\label{sec:discussion}

\para{Why the conjecture fails.} The conjecture in A380680 is word for word the conjecture of the companion entry
A380060, which was contributed on the same day and in which the conjecture is true (see the companion
paper~\cite{repo}). Replacing the series $A$ by the constant $2$ inside the sum looks like a small change, but
modulo $2$ it is not: the factor $(U-x^j)/(1-x^jU)$ is $1$ when $U\equiv1\pmod2$, and $x^j$ when $U\equiv0\pmod2$.
The published terms contradict the conjecture from $n=1$ on.

\para{Scope.} Theorem~\ref{thm:main} gives the remainder of every coefficient modulo $4$. We have not studied higher
moduli, and the theorem says nothing about the size of $a(n)$.

\para{A general principle.} The proof above, and the proof for the companion entry, are instances of one
observation. Reduction modulo $2$ is compatible with
all the operations in $\Sigma(U)$, so $\Sigma(U)$ modulo $2$ depends only on $U$ modulo $2$. By
Lemma~\ref{lem:double}, \emph{the remainder of $1+2\Sigma(U)$ modulo $4$ depends only on the remainder of $U$ modulo
$2$}. In particular, for every integer series $U$:
\begin{itemize}
\item if $U\equiv1\pmod2$ (odd constant term, all other coefficients even), then $1+2\Sigma(U)\equiv\theta(x)\pmod4$,
because $Q_j(1)=1$;
\item if $U\equiv0\pmod2$ (all coefficients even), then $1+2\Sigma(U)\equiv\theta(x^2)\pmod4$, because
$Q_j(0)=-x^j$ and $x^{j^2}(-x^j)^j=\pm x^{2j^2}$.
\end{itemize}
A380060 is the first case with $U=A$ and A380680 is the second case with $U=2$. We tested both cases numerically on
a few other choices of $U$; this general remark is not part of the Lean formalization.

\section*{About this paper}\label{sec:about}
\para{How the work was done.}
This paper is one of 28 papers on conjectures from the OEIS that came out of an AI research effort directed by
Dev~D.~Goyal~\cite{repo}. The roles were as follows.
\begin{itemize}
\item \emph{Dev~D.~Goyal} started the project, chose which problems to pursue, asked for proofs, formal verification and
revisions, and decided what to publish. Dev~D.~Goyal did not supply proof strategies or intermediate arguments.
\item \emph{ChatGPT} (OpenAI), using the GPT-6 model with Pro reasoning, worked autonomously from those prompts. It
searched the OEIS for suitable conjectures, found the proofs, wrote the Lean~4 formalization, and wrote a first draft
of this paper.
\item \emph{Claude Opus~5.5}, used with Claude Code, checked the work independently. It rebuilt the Lean project from source and confirmed that the theorems listed in Section~\ref{sec:checks} rely only on
Lean's standard axioms; compared the Lean definitions with the OEIS definition; recomputed the sequence with separately
written PARI/GP code, compared it with the terms published in the OEIS, and tested every claim of the paper on many
terms; and read every step of the proofs. It then wrote the present version of the paper, expanding the
explanations and, in places, writing out an argument in more detail or by a more elementary route. The theorems of
Section~\ref{sec:results} and the proof strategies are ChatGPT's. The Discussion section, including the short arguments
in it, was written by Claude Opus~5.5; its observations were checked by computer and are not part of the Lean
formalization.
\end{itemize}
Dev~D.~Goyal has reviewed the proofs and is responsible for this paper; this review is not an independent expert
validation.

\para{Review status.} This paper has not been peer reviewed, and it has not been submitted to a journal or to the OEIS.

\para{Priority.} An OEIS entry that labels a statement as a conjecture shows only that its contributor did not know
a proof; it does not show that no proof existed. During the research only limited searches for earlier work were
made. No claim of priority is made.

\para{Corrections.} Please report any error in the issue tracker of the repository~\cite{repo}, where the Lean sources and the numerical checks are also
available.

\section*{Acknowledgements}
The sequence and the conjecture are due to Paul~D.~Hanna~\cite{oeis}, who also contributed the companion entry
A380060~\cite{oeisother}.

\begin{thebibliography}{99}\small
\raggedright
\bibitem{oeis} P.~D.~Hanna, entry A380680 (G.f.\ $A(x)=\sum_{n=-\infty}^{\infty}x^{n^2}(2-x^n)^n/(1-2x^n)^n$), \emph{The On-Line Encyclopedia of Integer Sequences}, February 2025. \url{https://oeis.org/A380680}. Text as of October 2026.
\bibitem{oeisother} P.~D.~Hanna, entry A380060 (G.f.\ $A(x)$ satisfies $A(x)=\sum_{n=-\infty}^{\infty}x^{n^2}(A(x)-x^n)^n/(1-x^nA(x))^n$), \emph{The On-Line Encyclopedia of Integer Sequences}, February 2025. \url{https://oeis.org/A380060}. Text as of October 2026.
\bibitem{grosswald} E.~Grosswald, \emph{Representations of Integers as Sums of Squares}, Springer-Verlag, New York, 1985.
\bibitem{oeisdb} N.~J.~A.~Sloane and The OEIS Foundation Inc., \emph{The On-Line Encyclopedia of Integer Sequences}. \url{https://oeis.org}.
\bibitem{stanley} R.~P.~Stanley, \emph{Enumerative Combinatorics}, Volume~2, Cambridge University Press, 1999. (Chapter~5 treats formal power series, composition and compositional inverses.)
\bibitem{niven} I.~Niven, Formal power series, \emph{American Mathematical Monthly} 76 (1969), 871--889.
\bibitem{lean} L.~de~Moura and S.~Ullrich, The Lean~4 theorem prover and programming language, in \emph{Automated Deduction -- CADE 28}, Lecture Notes in Computer Science 12699, Springer, 2021, 625--635.
\bibitem{mathlib} The mathlib Community, The Lean mathematical library, in \emph{Proceedings of the 9th ACM SIGPLAN International Conference on Certified Programs and Proofs (CPP 2020)}, 367--381. Version used: commit \texttt{56a005855ea1}.
\bibitem{repo} D.~D.~Goyal, \emph{Proofs and corrections of OEIS conjectures}, repository, 2026. \url{https://github.com/dvgyl/oeis-proofs}.
\end{thebibliography}
\end{document}
